% Figura: construcciones para las condiciones de contorno (cilindro, circuito, camino del potencial).
\begin{figure}[ht]
\centering
\tikzset{medio/.style={font=\small, circle, draw=grispizarra, line width=0.4pt, inner sep=1.2pt, fill=white},
         lamina/.style={fill=grisclaro!60, draw=grispizarra!70, line width=0.4pt}}
% (a) cilindro de altura h
\begin{minipage}[t]{0.98\textwidth}\centering
\begin{tikzpicture}[baseline=(current bounding box.north), x=1cm, y=1cm]
  % vista oblicua
  \fill[lamina] (0,0) -- (3.4,0) -- (4.3,0.9) -- (0.9,0.9) -- cycle;
  \draw[vect, azulprof] (0.55,0.45) -- ++(0,1.5) node[above, font=\small] {$\un$};
  \node[medio] at (0.3,0.12) {2}; \node[medio] at (0.3,-0.35) {1};
  \draw[rojoteo, line width=0.7pt] (2.0,0.2) -- (2.0,0.95); \draw[rojoteo, line width=0.7pt] (2.9,0.2) -- (2.9,0.95);
  \draw[rojoteo, line width=0.7pt, fill=white, fill opacity=0.6] (2.45,0.95) ellipse (0.45 and 0.15);
  \draw[rojoteo, line width=0.7pt] (2.0,0.2) arc (180:360:0.45 and 0.15);
  \draw[rojoteo, line width=0.4pt, densely dashed] (2.0,0.2) arc (180:0:0.45 and 0.15);
  \node[font=\scriptsize, text=rojoteo, anchor=west] at (3.0,0.6) {$h$};
  % corte transversal
  \begin{scope}[shift={(6.3,0.2)}]
    \fill[lamina] (-0.2,-0.04) rectangle (3.6,0.04);
    \draw[grispizarra, line width=0.5pt] (-0.2,0) -- (3.6,0);
    \draw[azulprof, ->, line width=0.8pt] (-0.1,-0.5) -- (-0.1,1.0) node[above, font=\small] {$\un$};
    \node[medio] at (-0.45,0.35) {2}; \node[medio] at (-0.45,-0.35) {1};
    \draw[rojoteo, line width=0.9pt] (1.2,-0.6) rectangle (2.4,0.6);
    \draw[->, line width=0.9pt, azulacero] (1.8,0.6) -- ++(0,0.55) node[right, font=\scriptsize] {$\un_2=\un$};
    \draw[->, line width=0.9pt, azulacero] (1.8,-0.6) -- ++(0,-0.55) node[right, font=\scriptsize] {$\un_1=-\un$};
    \draw[<->, rojoteo, line width=0.5pt] (2.7,-0.6) -- (2.7,0.6) node[midway, right, font=\scriptsize] {$h$};
    \draw[<->, rojoteo, line width=0.5pt] (1.2,-0.85) -- (2.4,-0.85) node[midway, below, font=\scriptsize] {$\Delta S$};
  \end{scope}
\end{tikzpicture}\\[2pt]{\small\scshape (a)\; componente normal: cilindro de altura $h$ en la capa de transición}
\end{minipage}\\[10pt]
% (b) circuito
\begin{minipage}[t]{0.98\textwidth}\centering
\begin{tikzpicture}[baseline=(current bounding box.north), x=1cm, y=1cm]
  \fill[lamina] (0,0) -- (3.4,0) -- (4.3,0.9) -- (0.9,0.9) -- cycle;
  \draw[azulprof, ->, line width=0.8pt] (0.55,0.45) -- ++(0,1.5) node[above, font=\small] {$\un$};
  \node[medio] at (0.3,0.12) {2}; \node[medio] at (0.3,-0.35) {1};
  \draw[rojoteo, line width=0.8pt] (1.8,-0.3) rectangle (3.3,1.1);
  \begin{scope}[shift={(6.3,0.2)}]
    \fill[lamina] (-0.2,-0.04) rectangle (3.6,0.04);
    \draw[azulprof, ->, line width=0.8pt] (-0.1,-0.5) -- (-0.1,1.0) node[above, font=\small] {$\un$};
    \node[medio] at (-0.45,0.35) {2}; \node[medio] at (-0.45,-0.35) {1};
    \draw[rojoteo, line width=0.9pt, orienta] (1.0,-0.6) -- (1.0,0.6) -- (2.6,0.6) -- (2.6,-0.6) -- cycle;
    \draw[->, azulacero, line width=0.9pt] (1.8,0.6) -- ++(0.7,0) node[above, font=\scriptsize] {$\vc{t}$};
    \draw[->, azulacero, line width=0.9pt] (1.8,-0.6) -- ++(-0.7,0);
    \node[font=\scriptsize, text=azulacero, anchor=north] at (1.1,-0.64) {$\Delta\vc{l}_1=-\Delta l\,\vc{t}$};
    \node[font=\scriptsize, text=azulacero, anchor=south west] at (1.4,0.66) {$\Delta\vc{l}_2=\Delta l\,\vc{t}$};
    \node[font=\scriptsize, text=rojoteo, anchor=west] at (2.65,0) {$h$};
    \draw[azulprof, line width=0.7pt] (1.8,0) circle (0.12);
    \fill[azulprof] (1.8,0) circle (0.9pt);
    \node[font=\scriptsize, text=azulprof, anchor=north west] at (1.85,-0.08) {$\un'$};
  \end{scope}
\end{tikzpicture}\\[2pt]{\small\scshape (b)\; componente tangencial: circuito rectangular $C$ ($\un'$ sale de la figura)}
\end{minipage}\\[10pt]
% (c) camino del potencial
\begin{minipage}[t]{0.98\textwidth}\centering
\begin{tikzpicture}[baseline=(current bounding box.north), x=1cm, y=1cm]
  \fill[lamina] (0,0) -- (3.4,0) -- (4.3,0.9) -- (0.9,0.9) -- cycle;
  \draw[azulprof, ->, line width=0.8pt] (0.55,0.45) -- ++(0,1.5) node[above, font=\small] {$\un$};
  \node[medio] at (0.3,0.12) {2}; \node[medio] at (0.3,-0.35) {1};
  \draw[rojoteo, line width=0.9pt] (2.2,-0.55) -- (2.2,0.55);
  \fill[rojoteo] (2.2,0.55) circle (1.6pt) node[above right, font=\small] {$B$};
  \fill[rojoteo] (2.2,-0.55) circle (1.6pt) node[below right, font=\small] {$A$};
  \begin{scope}[shift={(6.3,0.2)}]
    \fill[lamina] (-0.2,-0.04) rectangle (3.6,0.04);
    \draw[azulprof, ->, line width=0.8pt] (-0.1,-0.5) -- (-0.1,1.0) node[above, font=\small] {$\un$};
    \node[medio] at (-0.45,0.35) {2}; \node[medio] at (-0.45,-0.35) {1};
    \draw[rojoteo, line width=0.9pt, orienta] (1.8,-0.6) -- (1.8,0.6);
    \fill[rojoteo] (1.8,0.6) circle (1.6pt) node[above right, font=\small] {$B$};
    \fill[rojoteo] (1.8,-0.6) circle (1.6pt) node[below right, font=\small] {$A$};
    \draw[<->, rojoteo, line width=0.5pt] (2.5,-0.6) -- (2.5,0.6) node[midway, right, font=\scriptsize] {$h$};
  \end{scope}
\end{tikzpicture}\\[2pt]{\small\scshape (c)\; potencial: camino $A\to B$ perpendicular a la superficie}
\end{minipage}
\caption{Construcciones geométricas para obtener las condiciones de contorno (la región 1 está
debajo de la superficie y la 2 encima; $\un$ apunta de 1 a 2). (a)~Un cilindro de bases de área $\Delta S$ paralelas a la
superficie y altura $h$, contenido en la capa de transición, para integrar la divergencia. (b)~Un
rectángulo, con dos lados de longitud $\Delta l$ paralelos a la superficie y dos de longitud $h$ (los
dos lados verticales), para integrar el rotacional. (c)~Un camino perpendicular a la
superficie, de $A$ a $B$, para el potencial.}
\label{fig:contorno}
\end{figure}
